\textbf{Время:} 00:30

\textbf{Тезис:} по импульсным корреляциям можно исследовать размеры источника частиц.

\textbf{На экране:} название → расходящиеся импульсы справа → три темы подзаголовка.

\textbf{Текст:} Я Балинов Темир. Сегодня поговорим о том, как перейти от измеренных импульсов частиц к размерам области, из которой они были испущены. На рисунке справа показаны именно импульсы. Нам понадобятся квантовая интерференция, статистика пар и моделирование в SMASH — эти три части перечислены в подзаголовке.

\textbf{Переход:} Сначала посмотрим, какой источник возникает при столкновении тяжёлых ядер.

\textbf{Справка:} Титульный рисунок — собственная тематическая схема, а не изображение зарегистрированного события.
